\documentclass[11pt]{article}

\usepackage{amsmath,amssymb,amsthm,mathtools}
\usepackage[margin=1in]{geometry}
\usepackage{microtype}

\newtheorem{theorem}{Theorem}[section]
\newtheorem{lemma}[theorem]{Lemma}
\newtheorem{corollary}[theorem]{Corollary}
\newtheorem{remark}[theorem]{Remark}

\newcommand{\dd}{\,\mathrm d}

\title{$K_4$ with a Two-Edge Tail:\\A Rooted Supporting-Plane Proof}
\author{}
\date{}

\begin{document}

\maketitle

\begin{abstract}
Let $H$ be obtained from $K_4$ by attaching an otherwise disjoint path of
length two at one clique vertex.  We prove
\[
 t(H,W)\geq p^3(2p-1)(3p-2)
\]
for every graphon of edge density $p\in[2/3,1]$, exactly the chromatic
target on its complete required interval.  Conditioning at the tail-bearing
clique vertex, applying the triangle Goodman bound inside its link, and using
the pointwise estimate $\tau\geq(2T_Wd-p)_+$ reduce the problem to a scalar
function of $d$ and $A=T_Wd$.  A sharp supporting plane at
$(d,A)=(p,p^2)$ integrates exactly because
$\int A=\int d^2$.  The scalar certificate is proved on the full feasible
region; its only lengthy sign is supplied by an explicit positive Bernstein
expansion.  This proves the formerly open NetworkX Atlas graph $142$ without
assuming the weighted link inequality that was left unresolved in the
failed-attempt ledger.
\end{abstract}


%%%%%%%%%%%%%%%%%%%%%%%%%%%%%%%%%%%%%%%%%%%%%%%%%%%%%%%%%%%%%%%%%%%%%%%%%%%%%%%
\section{The graph and its exact target}
%%%%%%%%%%%%%%%%%%%%%%%%%%%%%%%%%%%%%%%%%%%%%%%%%%%%%%%%%%%%%%%%%%%%%%%%%%%%%%%

Let $(\Omega,\mu)$ be an arbitrary probability space and let
$W\colon\Omega^2\to[0,1]$ be symmetric and measurable.  Write
\[
 p=\int_{\Omega^2}W(x,y)\dd\mu(x)\dd\mu(y),
 \qquad
 d(x)=\int_\Omega W(x,y)\dd\mu(y),
\]
and put
\[
 A(x)=(T_Wd)(x)=\int_\Omega W(x,y)d(y)\dd\mu(y).
\]
The rooted triangle density is
\[
 \tau(x)=\int_{\Omega^2}
 W(x,y)W(x,z)W(y,z)\dd\mu(y)\dd\mu(z).
\]

Let $H$ have clique vertices $o,a,b,c$ and tail vertices $u,v$, with
all six edges of the $K_4$ on $o,a,b,c$ and the two further edges
$ou,uv$.  After colouring the $K_4$, each tail vertex has $x-1$ choices.
Thus
\begin{equation}
 \chi_H(x)=x(x-1)^3(x-2)(x-3),
 \qquad \chi(H)=4.
 \label{eq:chromatic}
\end{equation}
For $q=1-p>0$, direct substitution gives
\begin{equation}
 \Phi_H(p)=q^6\chi_H(1/q)
 =p^3(2p-1)(3p-2).
 \label{eq:target}
\end{equation}
The expression on the right is its polynomial continuation at $p=1$.
The complete required interval is $p\in[2/3,1]$.


%%%%%%%%%%%%%%%%%%%%%%%%%%%%%%%%%%%%%%%%%%%%%%%%%%%%%%%%%%%%%%%%%%%%%%%%%%%%%%%
\section{The exact link reduction}
%%%%%%%%%%%%%%%%%%%%%%%%%%%%%%%%%%%%%%%%%%%%%%%%%%%%%%%%%%%%%%%%%%%%%%%%%%%%%%%

For $d(x)>0$, define the probability measure
\[
 \dd\mu_x(y)=\frac{W(x,y)}{d(x)}\dd\mu(y)
\]
and let $W_x$ denote the same kernel $W$ on this probability space.  Its
edge density is
\begin{equation}
 z_x=t(K_2,W_x)=\frac{\tau(x)}{d(x)^2}.
 \label{eq:link-edge}
\end{equation}
If $d(x)=0$, every homomorphism integrand having an edge incident with $x$
vanishes almost everywhere, so that point may simply be assigned zero.

Put, for every $z\geq0$,
\[
 f(z)=z(2z-1)_+.
\]
The triangle Goodman inequality, combined with nonnegativity, says that
every graphon $V$ of edge density $z$ satisfies
\begin{equation}
 t(K_3,V)\geq f(z).
 \label{eq:goodman}
\end{equation}
The function $f$ is nondecreasing: it is zero on $[0,1/2]$ and has
derivative $4z-1>0$ on $(1/2,\infty)$.

The elementary inequality $rs\geq r+s-1$ on $[0,1]^2$, applied to
$W(x,y),W(x,z)$ and multiplied by $W(y,z)$, gives
\begin{equation}
 \tau(x)\geq\max\{2A(x)-p,0\}.
 \label{eq:rooted-triangle}
\end{equation}
Indeed, the two positive terms integrate to $A(x)$ and the negative term
to $p$.  Also $\tau(x)\leq d(x)^2$, by dropping the factor $W(y,z)$.
Consequently the lower argument used in $f$ below never exceeds the actual
link density $z_x\leq1$.

Let $\kappa_4(x)$ be the rooted $K_4$ density at $o=x$.  Conditioning the
other three clique vertices through $\mu_x$ gives
\[
 \kappa_4(x)=d(x)^3t(K_3,W_x).
\]
Integrating the tail vertices first contributes $A(x)$, hence
\begin{equation}
 t(H,W)=\int_\Omega d(x)^3A(x)t(K_3,W_x)\dd\mu(x).
 \label{eq:link-identity}
\end{equation}
By \eqref{eq:goodman}, \eqref{eq:rooted-triangle}, and monotonicity of $f$,
\begin{equation}
 t(H,W)\geq\int_\Omega \mathcal F_p(d(x),A(x))\dd\mu(x),
 \label{eq:scalar-reduction}
\end{equation}
where
\begin{equation}
 \mathcal F_p(d,a)=
 \begin{cases}
 d^3a\,f\!\left(\dfrac{(2a-p)_+}{d^2}\right),&d>0,\\[6pt]
 0,&d=0.
 \end{cases}
 \label{eq:F}
\end{equation}
No pointwise division is made on the zero-degree set.

We shall use the following exact feasible information.
\begin{lemma}[Rooted feasible region]\label{lem:region}
For almost every $x$, with $d=d(x)$ and $a=A(x)$,
\[
 0\leq a\leq d,
 \qquad a\geq d+p-1.
\]
Moreover,
\[
 \int_\Omega d\dd\mu=p,
 \qquad
 \int_\Omega A\dd\mu=\int_\Omega d^2\dd\mu.
\]
\end{lemma}

\begin{proof}
Only the second pointwise inequality needs comment.  Applying
$rs\geq r+s-1$ to $W(x,y),d(y)$ and integrating gives
\[
 A(x)\geq\int_\Omega\bigl(W(x,y)+d(y)-1\bigr)\dd\mu(y)
 =d(x)+p-1.
\]
The upper bound follows from $d(y)\leq1$.  Finally, symmetry and Tonelli give
\[
 \int A(x)\dd\mu(x)
 =\int W(x,y)d(y)\dd\mu^2
 =\int d(y)^2\dd\mu(y).
\]
The first moment identity is the definition of $p$.
\end{proof}


%%%%%%%%%%%%%%%%%%%%%%%%%%%%%%%%%%%%%%%%%%%%%%%%%%%%%%%%%%%%%%%%%%%%%%%%%%%%%%%
\section{The sharp scalar supporting plane}
%%%%%%%%%%%%%%%%%%%%%%%%%%%%%%%%%%%%%%%%%%%%%%%%%%%%%%%%%%%%%%%%%%%%%%%%%%%%%%%

Define
\begin{align}
 T_p&=p^3(2p-1)(3p-2),\label{eq:T}\\
 \gamma_p&=p(20p^2-15p+2),\label{eq:gamma}\\
 \beta_p&=p^2(30p^2-21p+2),\label{eq:beta}
\end{align}
and
\begin{equation}
 L_p(d,a)=T_p+\beta_p(d-p)+\gamma_p(a-d^2).
 \label{eq:L}
\end{equation}
The coefficients are the derivatives, in the coordinates $(d,a-d^2)$,
of the active expression in \eqref{eq:F} at $(p,p^2)$.  In particular,
$L_p(p,p^2)=T_p$.

\begin{lemma}[Rooted $K_4$ supporting plane]\label{lem:scalar}
Let $p\in[2/3,1]$ and let $d,a\in[0,1]$ satisfy
\[
 0\leq a\leq d,
 \qquad a\geq d+p-1.
\]
Then
\begin{equation}
 \mathcal F_p(d,a)\geq L_p(d,a).
 \label{eq:support}
\end{equation}
\end{lemma}

\begin{proof}
First note that $\gamma_p>0$ on the stated interval.  When $d=0$,
feasibility forces $a=0$, and
\[
 L_p(0,0)=-2p^4(12p-7)\leq0=\mathcal F_p(0,0).
\]
Assume henceforth that $d>0$ and put
\[
 a_0=\frac p2+\frac{d^2}{4}.
\]
The scalar function is zero for $a\leq a_0$.  At the switching point,
direct expansion gives
\begin{equation}
 d\bigl(\mathcal F_p(d,a_0)-L_p(d,a_0)\bigr)
 =\frac{dp}{4}S_p(d),
 \label{eq:switch}
\end{equation}
where
\begin{align*}
 S_p(d)={}&3(20p^2-15p+2)d^2
 +(-120p^3+84p^2-8p)d\\
 &+96p^4-96p^3+30p^2-4p.
\end{align*}
Its leading coefficient is positive, and its discriminant in $d$ is
\[
 -8p(3p-2)^2(120p^3-120p^2+34p-3)\leq0.
\]
Indeed, if $u=p-2/3\geq0$, the final factor is
\[
 120u^3+120u^2+34u+\frac{17}{9}>0.
\]
Thus $S_p(d)\geq0$.  Since $L_p$ is increasing in $a$, for $a\leq a_0$
we obtain
\[
 L_p(d,a)\leq L_p(d,a_0)\leq0=\mathcal F_p(d,a),
\]
which proves the inactive case.

It remains to take $a\geq a_0$.  On this region, put
\begin{equation}
 G_{p,d}(a)=d\bigl(\mathcal F_p(d,a)-L_p(d,a)\bigr).
 \label{eq:G}
\end{equation}
Using $f(z)=z(2z-1)$ here gives the polynomial identity
\[
 d\mathcal F_p(d,a)=2a(2a-p)^2-d^2a(2a-p),
\]
so $G_{p,d}$ is a cubic in $a$ with leading coefficient $8$.
Its discriminant factors exactly as
\begin{equation}
 \operatorname{disc}_aG_{p,d}
 =-4dp(d-p)^2K(p,d).
 \label{eq:discriminant}
\end{equation}
We next give a finite exact positivity certificate for the remaining
factor.

For $0\leq u,d\leq1$, write
\[
 B_i^n(s)=\binom ni s^i(1-s)^{n-i}.
\]
With $u=3p-2$, exact coefficient conversion gives
\begin{equation}
 K\!\left(\frac{u+2}{3},d\right)
 =\frac1{765450}\sum_{i=0}^7\sum_{j=0}^6
 M_{ij}B_i^7(u)B_j^6(d),
 \label{eq:bernstein}
\end{equation}
where the rows of the positive integer matrix $M$ are
\[
\begin{array}{rrrrrrr}
9676800&21504000&24980480&26919060&33180000&48283550&69579300\\
18835200&50006400&60092480&65235210&79989840&116769575&176444850\\
33465600&109724800&136834880&149979960&181905520&262135400&403293600\\
56168640&230669280&299354832&333139662&399385152&560819025&856861470\\
90512640&469488960&635029632&720982620&856218816&1163172690&1733443740\\
141264000&931953600&1314243360&1527249330&1803325680&2366134875&3399784650\\
214617600&1813752000&2664912960&3175640640&3744513360&4758644700&6555605400\\
318427200&3474122400&5310794160&6495149430&7682260320&9512374725&12564861750
\end{array}
\]
Every Bernstein basis function is nonnegative, and in each coordinate the
basis functions sum to one.  Hence \eqref{eq:bernstein} proves
$K(p,d)>0$ throughout $p\in[2/3,1]$, $d\in[0,1]$.

If $d\ne p$, \eqref{eq:discriminant} is strictly negative.  The cubic
$G_{p,d}$ therefore has exactly one real root.  Equation \eqref{eq:switch}
is strict in this case: for $p>2/3$ the quadratic $S_p$ has negative
discriminant, while at $p=2/3$ one has
\[
 S_{2/3}(d)=\frac8{27}(3d-2)^2,
\]
whose only zero is $d=p$.  Thus $G_{p,d}(a_0)>0$.  A real cubic with positive
leading coefficient and exactly one real root is positive precisely to the
right of that root.  It follows that $G_{p,d}(a)>0$ for every $a\geq a_0$.

If $d=p$, direct factorization instead gives
\[
 G_{p,p}(a)=2(a-p^2)^2(4a+7p^2-4p).
\]
On $a\geq a_0=p/2+p^2/4$, the last factor is at least
$2p(4p-1)>0$.  Hence this expression is nonnegative as well.  Dividing by
$d>0$ completes the active case and the proof.
\end{proof}

\begin{remark}[Why this repairs, rather than assumes, the old missing step]
The earlier link approach sought the weighted mean inequality
\[
 \int dA\tau\geq\frac{2p-1}{p}\int d^3A.
\]
That statement does not follow from the accepted degree-weighted lemma,
because $A=T_Wd$ has no proved monotone ordering along complement edges.
Lemma~\ref{lem:scalar} never uses that comparison.  It keeps the pointwise
pair $(d,A)$, applies Goodman before averaging, and uses the exact
zero-integral correction $A-d^2$ in its supporting plane.
\end{remark}


%%%%%%%%%%%%%%%%%%%%%%%%%%%%%%%%%%%%%%%%%%%%%%%%%%%%%%%%%%%%%%%%%%%%%%%%%%%%%%%
\section{The graphon theorem and Atlas 142}
%%%%%%%%%%%%%%%%%%%%%%%%%%%%%%%%%%%%%%%%%%%%%%%%%%%%%%%%%%%%%%%%%%%%%%%%%%%%%%%

\begin{theorem}[$K_4$ with a two-edge tail]\label{thm:main}
Every graphon $W$ of edge density $p\in[2/3,1]$ satisfies
\[
 \boxed{
 t(H,W)\geq p^3(2p-1)(3p-2)=\Phi_H(p).
 }
\]
\end{theorem}

\begin{proof}
Combine \eqref{eq:scalar-reduction} with Lemmas~\ref{lem:region} and
\ref{lem:scalar}, and then integrate the supporting plane:
\begin{align*}
 t(H,W)
 &\geq\int\mathcal F_p(d,A)\\
 &\geq T_p+\beta_p\left(\int d-p\right)
       +\gamma_p\left(\int A-\int d^2\right)\\
 &=T_p.
\end{align*}
This is exactly \eqref{eq:target}.
\end{proof}

At $p=2/3$, the target is zero and the inequality is also immediate from
nonnegativity.  At $p=1$, $W=1$ almost everywhere and both sides equal one;
the proof above remains valid, but this observation handles the endpoint
without substituting into $1/(1-p)$.  Every balanced complete $k$-partite
graphon, $k\geq3$, realizes equality.  There $d=p$, $A=p^2$, the link is the
balanced complete $(k-1)$-partite graphon, and every link and scalar bound
used above is sharp.

\begin{corollary}[Atlas 142]
NetworkX Atlas graph $142$, with graph6 code \texttt{E\char126 AG}, is
positive.  It has order six, size eight, chromatic number four, edge list
\[
 01,\ 02,\ 03,\ 05,\ 12,\ 13,\ 23,\ 45,
\]
and chromatic polynomial
\[
 x(x-1)^3(x-2)(x-3).
\]
\end{corollary}

\begin{proof}
Vertices $0,1,2,3$ induce $K_4$, and $0-5-4$ is an otherwise disjoint
two-edge tail.  Thus this is exactly $H$.
\end{proof}

All integrands in the proof are bounded, nonnegative, and measurable, so
Tonelli--Fubini applies on an arbitrary probability space.  No atomlessness,
regularity, finite-step reduction, injective-copy interpretation, minimum
degree, or pointwise lower degree is used.  The only conditional quotient is
on $\{d>0\}$, and the zero-degree set is handled separately.

The theorem is specific to one two-edge tail attached at a $K_4$ vertex.  It
does not prove a general clique-tail family, multiple tails, arbitrary
attachments to cliques, or the discarded weighted link inequality above.
For formal verification, the reusable analytic unit is
Lemma~\ref{lem:scalar}; the graph-specific module awaiting inspection is
\texttt{lean/Taeyoung/Examples/Graph142.lean}.

The chromatic identity, rooted algebra, scalar factorizations, discriminant,
Bernstein conversion, and Atlas identification are reconstructed exactly by
\begin{verbatim}
python codes/verify_k4_two_edge_tail.py
\end{verbatim}


%%%%%%%%%%%%%%%%%%%%%%%%%%%%%%%%%%%%%%%%%%%%%%%%%%%%%%%%%%%%%%%%%%%%%%%%%%%%%%%
\section{What the Lean formalization actually proves}
%%%%%%%%%%%%%%%%%%%%%%%%%%%%%%%%%%%%%%%%%%%%%%%%%%%%%%%%%%%%%%%%%%%%%%%%%%%%%%%

The machine-checked development is
\texttt{lean/Taeyoung/Methods/K4Tail/\{Scalar,Link,Rows\}.lean}, ending at
\texttt{satisfiesLowerBound\_142}.  Theorem~\ref{thm:main} is proved in full, on
the whole interval $p\in[2/3,1]$, with no extra hypothesis.  The departure is
entirely inside Lemma~\ref{lem:scalar}.

\paragraph{The departure: neither the cubic discriminant nor the Bernstein
matrix is formalized.}
This note proves Lemma~\ref{lem:scalar} by treating
$G_{p,d}(a)=d(\mathcal F_p(d,a)-L_p(d,a))$ as a cubic in $a$, factoring
$\operatorname{disc}_aG_{p,d}=-4dp(d-p)^2K(p,d)$, certifying $K>0$ through the
degree-$(7,6)$ Bernstein expansion \eqref{eq:bernstein} with its $8\times7$
integer matrix, and then invoking the fact that a real cubic with positive
leading coefficient and exactly one real root is positive to the right of that
root.  None of those three ingredients appears in the formalization: not the
discriminant, not the matrix $M$, and not the root-counting argument, which has
no convenient Mathlib interface.

Instead the cubic is factored \emph{explicitly}.  In the coordinate
\[
 w=4a-2p-d^2,\qquad\text{so that}\qquad d\,\mathcal F_p(d,a)=\tfrac12\,a\,w\,(w+d^2),
\]
one has the polynomial identity, checked by \texttt{ring},
\begin{equation}
\begin{aligned}
 (11p-4)^3\cdot8\cdot\bigl(a(2a-p)w-d\,L_p(d,a)\bigr)
 ={}&\bigl((11p-4)w-R\bigr)^2\bigl((11p-4)w+2R+2(11p-4)(p+d^2)\bigr)\\
 &+(11p-4)(d-p)^2\,\Lambda\,w
 +(d-p)^2\,\Gamma,
\end{aligned}
 \label{eq:lean-factorization}
\end{equation}
where $R=(11p-4)r$ and $r(p,d)=\bigl(3p^2-2p\bigr)+\sigma(p)(d-p)$ is the
tangent line, at $d=p$, of the double root of the cubic; the slope
$\sigma(p)=(6p^2-9p+2)/(11p-4)$ is forced, and is the reason for the factor
$11p-4$.  The two remainder coefficients are
\begin{align*}
 \Lambda&=\bigl((11p-4)d+23p^2-22p+4\bigr)^2+\bigl(1658p^4-1850p^3+647p^2-60p-4\bigr),\\
 \Gamma&=Ad^2+Bd+C,\qquad A=-2(11p-4)(6p^2-9p+2)^2\leq0 .
\end{align*}
$\Lambda\geq0$ is a square plus a quartic in $p$ alone.  For $\Gamma$, the
feasibility $a\leq d$ together with $w\geq0$ gives $d^2\leq4d-2p$, and since
$A\leq0$ this yields $\Gamma\geq(4A+B)d+C-2Ap$; the same feasibility gives
$d\geq p/2$, and $4A+B\geq0$ then leaves $\Gamma\geq\tfrac12Bp+C$, again a
polynomial in $p$ alone.

Every one-variable inequality that remains --- there are seven --- has
\emph{all} coefficients nonnegative after $p=\tfrac23+u$, so each is a one-line
\texttt{nlinarith}.  In particular the quantity
$1658p^4-1850p^3+647p^2-60p-4$, whose positivity replaces the discriminant
condition of this note, is $\tfrac{1856}{81}+\tfrac{8128}{27}u+\tfrac{4105}{3}u^2+\tfrac{7714}{3}u^3+1658u^4$.

The two routes prove the same lemma on the same interval.  The note's is the
more informative one --- it exhibits the root structure --- while
\eqref{eq:lean-factorization} is the one a proof assistant can check without a
theory of cubics.

\paragraph{The truncations are carried, not divided away.}
The note defines $\mathcal F_p(d,0)=0$ by hand and notes that no pointwise
division is made on the zero-degree set.  The formalization avoids the
definition entirely: \texttt{K4Tail.mul\_rootedK4\_ge} states the link bound in
the form
\[
 d(x)\,\kappa_4(x)\;\geq\;u\cdot(2u-d(x)^2)_+,\qquad u=(2A(x)-p)_+,
\]
multiplied through by $d(x)$, so both sides are polynomials in the truncated
quantities and $d(x)=0$ is an ordinary case rather than a convention.  The
monotonicity of $f$ used in \eqref{eq:scalar-reduction} becomes the elementary
step that $t\mapsto t(2t-d^2)$ is nondecreasing once $2t\geq d^2$, applied to
$\tau\geq u$.

\paragraph{Everything else transcribes directly.}
\eqref{eq:goodman} is \texttt{K4Tail.goodman}, the pure-chordal clique bound at
$r=3$ extended below the threshold by nonnegativity; \eqref{eq:link-identity} is
the existing cone identity \texttt{Link.rootedDensity\_eq} at $F=K_3$;
\eqref{eq:rooted-triangle} is \texttt{Link.rootedTriangle\_ge}; and the two
moment identities of Lemma~\ref{lem:region} are \texttt{Foundation.integral\_degree}
and \texttt{BookTail.integral\_pathOp}.  The final integration of the supporting
plane is \texttt{K4Tail.integral\_plane}, which is exactly the statement that the
two corrections $d-p$ and $A-d^2$ have mean zero.

\end{document}
